% sec-00: abstract, how to read, figure and table inventory.
% Resolved from: applications/README.md 1 and 2, science-golden-age chunk-01,
% daraxonrasib-llm-story.md, and the five diagram stage READMEs.
\section*{Abstract}
\phantomsection\addcontentsline{toc}{section}{Abstract}

In July 2026 the White House Office of Science and Technology Policy asked the
federal research portfolio, approximately \$200 billion annually, to prioritise
the individual scientist over legacy institutions
\cite{kratsios2026sciencegoldenage}. This paper reports what happens when one
unaffiliated scientist takes that instruction literally and writes ten funding
applications for the same Phase 1 pancreatic cancer trial, each to a different
recipient, each addressing the mechanism that recipient runs.

The trial is constant across all ten: an open-label, single-arm, 3+3 dose
escalation of perioperative daraxonrasib (RMC-6236) in up to 18 participants
with resectable or borderline-resectable KRAS G12-mutated pancreatic ductal
adenocarcinoma, performed as a staged eight-arm robotic pancreaticoduodenectomy
with an on-premises large language model confined to advisory output
\cite{onpremwhippl,phase1ind}. Five applications lead with the operation and
five with the drug and patient selection; both sets describe the same hybrid
procedure. UC San Diego Moores Cancer Center is named in all ten as the intended
partner of choice, approached at the feasibility stage only, and is described
nowhere as a partner, sponsor, trial site, or endorser.

The evidence base is stated with its limits attached. An August 2025 quantitative
systems pharmacology simulation returned median overall survival of 12.8 months
for a daraxonrasib combination against 5.4 months for chemotherapy
\cite{qspmetpancre}; in May 2026 the RASolute 302 trial reported 13.2 against
6.6 months in the RAS G12 metastatic population \cite{rasolute302}. That
proximity is a chronology observation, not a validation claim, and the three
differences are stated wherever it appears. Neither result speaks to the
resectable setting, which is what the proposed trial would measure.

The applications themselves carry no digital object identifiers. This paper
carries one.

\bmhead{How to Read This Paper}

This is a summary of ten application file sets, not an application. It reports
what was written, to whom, and on what evidence, and it is written so that a
reader can check every claim against a named file in a public repository
\cite{repo2026}.

It is not medical advice, not regulatory advice, and not a claim that any
institution, sponsor, or drug developer has agreed to anything. No patient has
been treated. The trial does not exist until a site, an institutional review
board, an active investigational new drug application, monitoring, insurance,
and a sponsor of record are all in place.

Three reading paths. A \textbf{funder} should read \S2 for the recipient set and
\S8 for the budget in its ten shapes. A \textbf{clinician} should read \S6 for
the governance boundaries and \S5 for the evidence with its limitations. A
\textbf{metascience reader} should read \S1 for the policy argument and \S9 for
the build method, which is reproducible from the prompt and the sources alone.

\bmhead{Figure and Table Inventory}

Twenty figures across five machine-readable diagram platforms. The split follows
what each figure has to answer rather than an equal quota: an even division
across five platforms would have forced at least four figures into a vocabulary
that cannot state their subject.

\begin{apptable}
\begin{tabularx}{\textwidth}{>{\raggedright\arraybackslash}p{0.9cm} >{\raggedright\arraybackslash}p{2.9cm} >{\raggedright\arraybackslash}p{1.1cm} >{\raggedright\arraybackslash}X}
\headrow \thead{Fig} & \thead{Type} & \thead{Sec} & \thead{The question it answers}\\
\midrule
1 & mermaid & \S1 & How does one policy sentence become ten addressed applications? \\
2 & plantuml & \S1 & Where does an unaffiliated proposal stall, and what removes each stall? \\
3 & d2 & \S2 & How do the ten compare on mechanism, term, ask, and lead? \\
4 & mermaid & \S5 & How long did each activity run, and what overlapped? \\
5 & plantuml & \S3 & Who is authorised to take which action? \\
6 & graphviz & \S2 & Which award unblocks which activity, and what is redundant? \\
7 & d2 & \S5 & Which evidence tier contains which, and what may not be cited as what? \\
8 & mermaid & \S3 & What are the reviewer's gates, and which section answers each? \\
9 & diagrams & \S6 & What runs inside the trust boundary, and what may cross it? \\
10 & graphviz & \S6 & By what path could an unintended motion reach the patient? \\
11 & d2 & \S8 & Where does each dollar land, and what is contributed rather than paid? \\
12 & mermaid & \S6 & Who says what to whom across one operative day? \\
13 & plantuml & \S6 & Under what guard does the advisory system change state? \\
14 & diagrams & \S6 & Which model produced which artifact, and who reviewed it? \\
15 & d2 & \S6 & What joins to what in the released dataset? \\
16 & graphviz & \S5 & Which prior work fed which, and where do the two lines meet? \\
17 & mermaid & \S8 & How long is each review clock against the site milestone? \\
18 & diagrams & \S9 & How does one operator's stack compare with a program office? \\
19 & plantuml & \S4 & What runs in parallel once any one award lands? \\
20 & graphviz & \S10 & Which field can falsify the programme's central claim? \\
\end{tabularx}
\end{apptable}
\tabcap{The twenty figures, their platform, their section, and the question each
exists to answer. Six are mermaid-type, four d2-type, four graphviz-type, three
plantuml-type, and three diagrams-python-type.}

\begin{apptable}
\begin{tabularx}{\textwidth}{>{\raggedright\arraybackslash}p{1.1cm} >{\raggedright\arraybackslash}p{1.1cm} >{\raggedright\arraybackslash}X >{\raggedright\arraybackslash}p{4.3cm}}
\headrow \thead{Table} & \thead{Sec} & \thead{Contents} & \thead{Source directory}\\
\midrule
1 to 2 & \S0 & Figure and table inventories & The five stage directories \\
3 to 4 & \S1 & Where the enterprise stalls; the four goals answered & \appfile{funding/science-golden-age/} \\
5 to 6 & \S2 & The ten recipients; what differs and what does not & \appfile{applications/} \\
7 & \S3 & Set A figure inventory by type & \appfile{applications/app-01} to \appfile{app-05} \\
8 & \S4 & Consortium measurements against resection timing & \appfile{applications/app-06} \\
9 to 11 & \S5 & Results with stated limitations; cost and schedule; the fourteen works & \appfile{funding/supplementary/source-files/} \\
12 & \S6 & The three frontier-model roles & \appfile{funding/tripartisan-llm-support.md} \\
13 to 14 & \S7 & Three positioning corrections; two site routes & \appfile{funding/potential-partners/} \\
15 & \S8 & One budget in ten shapes & \appfile{funding/RFA-RM-27-001-v2/} \\
16 to 17 & \S9 & The thirteen sub-prompts; the diagram split & \appfile{sub-prompts/} \\
18 to 19 & \S10 & What is not claimed; risks the programme cannot retire alone & Application back matter \\
\end{tabularx}
\end{apptable}
\tabcap{The paper's nineteen tables, the section each appears in, and the
repository directory its values are drawn from. No value in any table is
computed here for the first time.}
